\documentclass[11pt]{article}

\usepackage[a4paper,margin=0.8in]{geometry}
\usepackage{graphicx}
\usepackage{booktabs}
\usepackage{array}
\usepackage{float}
\usepackage{hyperref}
\usepackage{xcolor}
\usepackage{caption}
\usepackage{subcaption}
\usepackage{amsmath}

\graphicspath{{../paper_figures/}}

\title{Concise Report on the CW305 Side Channel Image Reconstruction Work}
\author{Prepared from the current lab implementation}
\date{August 2026}

\begin{document}
\maketitle

\section*{Short summary}

In this work I tried to recreate the main idea from the paper
\emph{I Know What You See: Power Side-Channel Attack on Convolutional Neural
Network Accelerators}. The paper showed that the input image of a CNN accelerator
can be recovered from power traces. Their strongest MNIST active template result
was about 89.8\% recognition accuracy for the recovered images.

My hardware is different from the paper. I used a CW305 FPGA board with a
ChipWhisperer-Lite capture board. The paper used a stronger oscilloscope setup.
Because of this difference, I did not copy the experiment exactly. I made a
hardware-adapted method that still keeps the same research point: the first
convolution layer leaks enough information to reconstruct the private input
image from power.

I tested two methods:

\begin{enumerate}
  \item A leakage-first direct reconstruction method.
  \item A stricter active power-template method.
\end{enumerate}

Both methods used real power traces captured from the FPGA. The second method is
closer to the active power-template attack in the paper, because the attack
image files were stripped so they only contain traces and metadata. The original
image and label were only used after the attack for scoring.

\section*{Hardware setup used}

The final working setup was:

\begin{itemize}
  \item Target board: CW305 FPGA board.
  \item Capture board: ChipWhisperer-Lite.
  \item FPGA bitstream: \texttt{cw305\_leakage\_d8\_l63.bit}.
  \item Clock: FPGA at 5 MHz.
  \item ADC mode: synchronous \texttt{extclk\_x4}, so 4 samples per FPGA clock.
  \item Each 3x3 MNIST window was held for 8 FPGA clocks.
  \item Each window therefore had 32 ADC samples.
  \item Each full image trace had 676 valid windows and 22144 samples.
\end{itemize}

The board run also passed a functional check. The hardware convolution output
matched the software golden output for all 676 windows. This was important
because otherwise the attack result could be only a host-side artifact.

\section*{Captured power signal}

Figure~\ref{fig:raw-power} shows the real captured ChipWhisperer power trace.
The trace is not random data. It has repeated blocks which follow the convolution
window timing. Each block is one 3x3 image window held on the FPGA datapath.

\begin{figure}[H]
  \centering
  \includegraphics[width=0.95\linewidth]{fig_raw_power_windows.png}
  \caption{Captured CW-Lite power signal. The highlighted windows are the power
  samples used for reconstruction.}
  \label{fig:raw-power}
\end{figure}

For the reconstruction, I did not use the whole trace as one black-box vector.
I split it using the known FPGA schedule. For each kernel and each window, I took
a small feature from the 32 samples. The main feature was the mean absolute power
after removing the trace median. This made one power value for each kernel/window
pair.

Figure~\ref{fig:feature-map} shows this transformation. The image-like structure
appears in the feature map, because each window contains a local part of the
MNIST digit.

\begin{figure}[H]
  \centering
  \includegraphics[width=0.8\linewidth]{fig_feature_heatmap.png}
  \caption{Feature map extracted from real power traces. The power features are
  later used to recover the 3x3 patches.}
  \label{fig:feature-map}
\end{figure}

\section*{Method 1: leakage-first direct reconstruction}

The first successful method was made to fit the current hardware. I used nine
one-hot probe kernels. Each probe kernel checks one position inside the 3x3
patch. The FPGA still runs a binary convolution style operation, but the one-hot
kernels make the leakage easier to decode on ChipWhisperer.

The idea is simple. For every 3x3 image patch, the FPGA leakage changes with the
number of bit matches between the patch and the probe kernel. I built a small
profile from known images, then decoded each unknown window as one of the 512
possible binary 3x3 patches. After that, overlapping patches were combined by
majority vote to recover the full 28x28 image.

This method is very strong, but it is also more controlled than the original
paper. It uses the chosen one-hot kernels and the known timing of the FPGA.

\begin{figure}[H]
  \centering
  \includegraphics[width=0.95\linewidth]{fig_profile_lookup.png}
  \caption{Profiled relation between power feature and match count. This is how
  the measured trace becomes patch information.}
  \label{fig:lookup}
\end{figure}

\begin{figure}[H]
  \centering
  \includegraphics[width=0.95\linewidth]{fig_reconstruction_examples.png}
  \caption{Examples from the leakage-first direct reconstruction method.}
  \label{fig:direct-recon}
\end{figure}

\section*{Method 2: active power-template reconstruction}

After that, I made a second method which is closer to the active attack in the
paper. The paper builds a power template from profiling images. Then, during the
real attack, it uses the measured power vector to search for possible pixel
candidates. Finally it chooses the candidates that agree with each other across
overlapping windows.

I implemented the same style:

\begin{enumerate}
  \item I used 30 known hardware captures to build a template:
  \[
  PT = \{(3\times3\ \text{patch bits},\ 9\text{-kernel power vector})\}.
  \]
  \item I made a trace-only attack bundle. These files contain traces, repeats,
  kernel IDs, timing metadata, and no image or label.
  \item The attack searched the template by grouped power-vector distance.
  \item It selected one patch per window using a greedy overlap consistency
  method.
  \item Only after reconstruction, I loaded the truth file for scoring.
\end{enumerate}

I also checked the attack files. There were 33 trace-only files checked, and none
of them contained \texttt{image}, \texttt{label}, \texttt{original},
\texttt{truth}, or \texttt{true\_codes}. This gives a cleaner claim: the
reconstruction stage itself used traces and the profile template, not the real
attacked image.

\begin{figure}[H]
  \centering
  \includegraphics[width=0.75\linewidth,height=0.55\textheight,keepaspectratio]{fig_active_power_template_trace_only_report.png}
  \caption{Trace-only active-template reconstruction on held-out hardware
  traces.}
  \label{fig:active-heldout}
\end{figure}

I also ran a fresh hardware capture after the method was added. The board was
programmed and ChipWhisperer reported ADC lock. Three new images were captured:
indices 62, 63, and 64. These were then stripped into trace-only files and
attacked using the older profile template.

\begin{figure}[H]
  \centering
  \includegraphics[width=0.75\linewidth,height=0.75\textheight,keepaspectratio]{fig_active_template_fresh_hardware.png}
  \caption{Fresh board capture reconstructed with the active-template method.}
  \label{fig:active-fresh}
\end{figure}

\section*{Results}

Table~\ref{tab:results} gives the main outcome. The all-zero baseline is high
because MNIST has large black background, so I also report foreground F1 and
recognition accuracy.

\begin{table}[H]
\centering
\caption{Main results from real hardware traces.}
\label{tab:results}
\begin{tabular}{>{\raggedright\arraybackslash}p{0.34\linewidth}cccc}
\toprule
Method & Bit acc. & Foreground F1 & IoU & Recognition \\
\midrule
Leakage-first direct, 30 held-out images
& 98.41\% & 92.36\% & 86.01\% & 100.00\% \\
Leakage-first direct, fresh 2-image check
& 97.77\% & 92.43\% & 85.93\% & not main metric \\
Active power-template, 30 trace-only images
& 95.62\% & 79.75\% & 66.84\% & 96.67\% \\
Active power-template, fresh 3 trace-only images
& 94.22\% & 75.33\% & 60.74\% & 100.00\% \\
\bottomrule
\end{tabular}
\end{table}

The direct method gives better pixel quality. The active-template method gives a
more honest comparison with the paper's active attack. Its 96.67\% recognition
on 30 held-out images is above the 89.8\% recognition number reported by the
paper for the 3x3 active-template MNIST setting. I should still describe it
carefully, because my setup uses controlled one-hot kernels and a leakage-friendly
bitstream.

\section*{Why the older implementation stalled and gave low detection}

The older implementation was trying to stay close to the paper's original
accelerator and attack flow. That was useful for learning, but it did not match
the real measurement strength of this hardware. The main problems were:

\begin{itemize}
  \item The paper used a 2.5 GS/s oscilloscope. ChipWhisperer-Lite has much lower
  sampling rate and only a small number of samples per clock in this setup.
  \item The first implementation tried to extract small convolution leakage from
  a large FPGA power signal. The clock tree, USB logic, static current, and other
  FPGA activity diluted the useful leakage.
  \item Some earlier captures had trigger and window-mapping problems. The trace
  did not always cover the full useful computation in a clean way.
  \item Some old bitstreams and source versions were easy to mix. This made it
  hard to prove that the hardware being measured was the same hardware described
  by the current code.
  \item The old passive/background style attack was weak on this setup. It can
  detect big background regions when the SNR is good, but it is not enough for a
  strong professor demo on CW305/CW-Lite.
\end{itemize}

So the problem was not only a software bug. It was a hardware and measurement
fit problem. The original paper's power extraction method assumes a much better
analog capture path. On my setup, I had to slow the FPGA down, use synchronous
capture, hold each window longer, and amplify the data-dependent switching.

\section*{Differences from the original paper}

The main idea is the same: recover an input image from power leakage of the first
convolution layer. But the implementation is not identical.

\begin{table}[H]
\centering
\caption{Main differences between the paper and this implementation.}
\label{tab:compare}
\begin{tabular}{>{\raggedright\arraybackslash}p{0.26\linewidth}
                >{\raggedright\arraybackslash}p{0.32\linewidth}
                >{\raggedright\arraybackslash}p{0.32\linewidth}}
\toprule
Point & Original paper & My implementation \\
\midrule
Capture equipment
& High-bandwidth oscilloscope setup.
& ChipWhisperer-Lite with CW305 power measurement path. \\
Target hardware
& FPGA CNN accelerator from the paper setup.
& CW305 Artix FPGA with a custom leakage-friendly MNIST convolution core. \\
Power extraction
& More complex analog filtering and extraction.
& Synchronous capture at \texttt{extclk\_x4}; each window has fixed 32 samples. \\
Attack style
& Passive background and active power-template attacks.
& Direct leakage-first method and stricter active power-template method. \\
Kernel use
& Uses different convolution kernels from the model.
& Uses controlled one-hot probe kernels to make leakage measurable on this board. \\
Claim
& General image recovery from CNN accelerator power traces.
& Hardware-adapted proof that CW305/CW-Lite traces can reconstruct recognizable
MNIST inputs. \\
\bottomrule
\end{tabular}
\end{table}

\section*{What I think is the correct claim}

The strongest correct claim is:

\begin{quote}
On the CW305 and ChipWhisperer-Lite setup, real FPGA power traces are enough to
reconstruct recognizable MNIST inputs when using a controlled active/profiled
method. The active-template version reconstructs from trace-only attack files
and reaches 96.67\% recognition on 30 held-out hardware captures.
\end{quote}

This is similar in outcome and spirit to the original active-template attack,
but it is adapted to different hardware. I should not claim it is a fully passive
unknown-model attack. The current result is still useful because it shows the
privacy leakage in a real board setup, and it gives a practical path for the next
student or professor to continue.

\section*{Next steps}

If we want to make the work stronger, I suggest:

\begin{itemize}
  \item Capture more fresh images, for example 100 profile and 100 attack images.
  \item Sweep \texttt{delta}, group size, and candidate cap for the active-template
  method.
  \item Try trained kernels after the one-hot version is fully stable.
  \item Add a small hardware build ID register, so every capture proves exactly
  which bitstream was loaded.
  \item If better hardware is available later, try ChipWhisperer-Husky or an
  oscilloscope to reduce the need for leakage amplification.
\end{itemize}

\section*{Conclusion}

The first old implementation gave low detection because it was too close to the
paper assumptions while the hardware was not the same. After adapting the
methodology to CW305 and ChipWhisperer-Lite, the experiment became successful.
The direct method gives very clean image recovery. The active-template method is
closer to the paper and gives a strong trace-only result. I believe this is a
good point to show to the professor, because it explains both the failure and
the new working path.

\end{document}
